% ৩.২১ কাউন্টার
\section{Counter}

\subsection{সংজ্ঞা}
\begin{itemize}
\item Pulse গণনা করে; Flip-flop দ্বারা নির্মিত sequential circuit
\end{itemize}

\subsection{প্রকারভেদ}
\begin{itemize}
\item Asynchronous, Synchronous, Up, Down, Mod, Up-Down
\end{itemize}

\subsection{ব্যবহার}
\begin{itemize}
\item Digital clock, frequency counter, timers
\end{itemize}

\begin{learningbox}
এই টপিক শেষে শিক্ষার্থী পারবে:
\begin{itemize}
\item counter-এর সংজ্ঞা ও sequential circuit ধারণা ব্যাখ্যা করতে।
\item asynchronous ও synchronous counter আলাদা করতে।
\item mod counter, up counter ও down counter-এর ধারণা লিখতে।
\end{itemize}
\end{learningbox}

\begin{definitionbox}{Counter}
clock pulse বা event গণনা করে নির্দিষ্ট binary sequence তৈরি করে যে sequential circuit, তাকে counter বলা হয়। counter সাধারণত flip-flop দিয়ে তৈরি হয়।
\end{definitionbox}

\begin{center}
\begin{tabular}{|p{0.24\textwidth}|p{0.30\textwidth}|p{0.30\textwidth}|}
\hline
\textbf{ধরন} & \textbf{বৈশিষ্ট্য} & \textbf{উদাহরণ} \\
\hline
Asynchronous & flip-flop একসাথে clock পায় না & ripple counter \\
\hline
Synchronous & সব flip-flop একই clock পায় & high-speed counter \\
\hline
Up counter & sequence বাড়ে & 000, 001, 010 ... \\
\hline
Down counter & sequence কমে & 111, 110, 101 ... \\
\hline
Mod counter & নির্দিষ্ট state পর্যন্ত গুনে repeat করে & Mod-10 decade counter \\
\hline
\end{tabular}
\end{center}

\begin{center}
\fbox{\parbox{0.88\textwidth}{
\centering
\textbf{State diagram placeholder}\\[4pt]
3-bit up counter: 000 $\rightarrow$ 001 $\rightarrow$ 010 $\rightarrow$ 011 $\rightarrow$ 100 $\rightarrow$ 101 $\rightarrow$ 110 $\rightarrow$ 111 $\rightarrow$ 000
}}
\end{center}

\begin{boardbox}{বোর্ড ফোকাস}
Counter প্রশ্নে ``flip-flop'', ``clock pulse'', ``sequential circuit'' এবং ``state sequence'' শব্দগুলো ব্যবহার করলে উত্তর বেশি স্পষ্ট হয়।
\end{boardbox}
